% Options for packages loaded elsewhere
\PassOptionsToPackage{unicode}{hyperref}
\PassOptionsToPackage{hyphens}{url}
\documentclass[
  11pt,
]{article}
\usepackage{xcolor}
\usepackage[left=3cm,right=3cm,top=2cm,bottom=2cm]{geometry}
\usepackage{amsmath,amssymb}
\setcounter{secnumdepth}{5}
\usepackage{iftex}
\ifPDFTeX
  \usepackage[T1]{fontenc}
  \usepackage[utf8]{inputenc}
  \usepackage{textcomp} % provide euro and other symbols
\else % if luatex or xetex
  \usepackage{unicode-math} % this also loads fontspec
  \defaultfontfeatures{Scale=MatchLowercase}
  \defaultfontfeatures[\rmfamily]{Ligatures=TeX,Scale=1}
\fi
\usepackage{lmodern}
\ifPDFTeX\else
  % xetex/luatex font selection
\fi
% Use upquote if available, for straight quotes in verbatim environments
\IfFileExists{upquote.sty}{\usepackage{upquote}}{}
\IfFileExists{microtype.sty}{% use microtype if available
  \usepackage[]{microtype}
  \UseMicrotypeSet[protrusion]{basicmath} % disable protrusion for tt fonts
}{}
\makeatletter
\@ifundefined{KOMAClassName}{% if non-KOMA class
  \IfFileExists{parskip.sty}{%
    \usepackage{parskip}
  }{% else
    \setlength{\parindent}{0pt}
    \setlength{\parskip}{6pt plus 2pt minus 1pt}}
}{% if KOMA class
  \KOMAoptions{parskip=half}}
\makeatother
\usepackage{color}
\usepackage{fancyvrb}
\newcommand{\VerbBar}{|}
\newcommand{\VERB}{\Verb[commandchars=\\\{\}]}
\DefineVerbatimEnvironment{Highlighting}{Verbatim}{commandchars=\\\{\}}
% Add ',fontsize=\small' for more characters per line
\usepackage{framed}
\definecolor{shadecolor}{RGB}{248,248,248}
\newenvironment{Shaded}{\begin{snugshade}}{\end{snugshade}}
\newcommand{\AlertTok}[1]{\textcolor[rgb]{0.94,0.16,0.16}{#1}}
\newcommand{\AnnotationTok}[1]{\textcolor[rgb]{0.56,0.35,0.01}{\textbf{\textit{#1}}}}
\newcommand{\AttributeTok}[1]{\textcolor[rgb]{0.13,0.29,0.53}{#1}}
\newcommand{\BaseNTok}[1]{\textcolor[rgb]{0.00,0.00,0.81}{#1}}
\newcommand{\BuiltInTok}[1]{#1}
\newcommand{\CharTok}[1]{\textcolor[rgb]{0.31,0.60,0.02}{#1}}
\newcommand{\CommentTok}[1]{\textcolor[rgb]{0.56,0.35,0.01}{\textit{#1}}}
\newcommand{\CommentVarTok}[1]{\textcolor[rgb]{0.56,0.35,0.01}{\textbf{\textit{#1}}}}
\newcommand{\ConstantTok}[1]{\textcolor[rgb]{0.56,0.35,0.01}{#1}}
\newcommand{\ControlFlowTok}[1]{\textcolor[rgb]{0.13,0.29,0.53}{\textbf{#1}}}
\newcommand{\DataTypeTok}[1]{\textcolor[rgb]{0.13,0.29,0.53}{#1}}
\newcommand{\DecValTok}[1]{\textcolor[rgb]{0.00,0.00,0.81}{#1}}
\newcommand{\DocumentationTok}[1]{\textcolor[rgb]{0.56,0.35,0.01}{\textbf{\textit{#1}}}}
\newcommand{\ErrorTok}[1]{\textcolor[rgb]{0.64,0.00,0.00}{\textbf{#1}}}
\newcommand{\ExtensionTok}[1]{#1}
\newcommand{\FloatTok}[1]{\textcolor[rgb]{0.00,0.00,0.81}{#1}}
\newcommand{\FunctionTok}[1]{\textcolor[rgb]{0.13,0.29,0.53}{\textbf{#1}}}
\newcommand{\ImportTok}[1]{#1}
\newcommand{\InformationTok}[1]{\textcolor[rgb]{0.56,0.35,0.01}{\textbf{\textit{#1}}}}
\newcommand{\KeywordTok}[1]{\textcolor[rgb]{0.13,0.29,0.53}{\textbf{#1}}}
\newcommand{\NormalTok}[1]{#1}
\newcommand{\OperatorTok}[1]{\textcolor[rgb]{0.81,0.36,0.00}{\textbf{#1}}}
\newcommand{\OtherTok}[1]{\textcolor[rgb]{0.56,0.35,0.01}{#1}}
\newcommand{\PreprocessorTok}[1]{\textcolor[rgb]{0.56,0.35,0.01}{\textit{#1}}}
\newcommand{\RegionMarkerTok}[1]{#1}
\newcommand{\SpecialCharTok}[1]{\textcolor[rgb]{0.81,0.36,0.00}{\textbf{#1}}}
\newcommand{\SpecialStringTok}[1]{\textcolor[rgb]{0.31,0.60,0.02}{#1}}
\newcommand{\StringTok}[1]{\textcolor[rgb]{0.31,0.60,0.02}{#1}}
\newcommand{\VariableTok}[1]{\textcolor[rgb]{0.00,0.00,0.00}{#1}}
\newcommand{\VerbatimStringTok}[1]{\textcolor[rgb]{0.31,0.60,0.02}{#1}}
\newcommand{\WarningTok}[1]{\textcolor[rgb]{0.56,0.35,0.01}{\textbf{\textit{#1}}}}
\usepackage{graphicx}
\makeatletter
\newsavebox\pandoc@box
\newcommand*\pandocbounded[1]{% scales image to fit in text height/width
  \sbox\pandoc@box{#1}%
  \Gscale@div\@tempa{\textheight}{\dimexpr\ht\pandoc@box+\dp\pandoc@box\relax}%
  \Gscale@div\@tempb{\linewidth}{\wd\pandoc@box}%
  \ifdim\@tempb\p@<\@tempa\p@\let\@tempa\@tempb\fi% select the smaller of both
  \ifdim\@tempa\p@<\p@\scalebox{\@tempa}{\usebox\pandoc@box}%
  \else\usebox{\pandoc@box}%
  \fi%
}
% Set default figure placement to htbp
\def\fps@figure{htbp}
\makeatother
% definitions for citeproc citations
\NewDocumentCommand\citeproctext{}{}
\NewDocumentCommand\citeproc{mm}{%
  \begingroup\def\citeproctext{#2}\cite{#1}\endgroup}
\makeatletter
 % allow citations to break across lines
 \let\@cite@ofmt\@firstofone
 % avoid brackets around text for \cite:
 \def\@biblabel#1{}
 \def\@cite#1#2{{#1\if@tempswa , #2\fi}}
\makeatother
\newlength{\cslhangindent}
\setlength{\cslhangindent}{1.5em}
\newlength{\csllabelwidth}
\setlength{\csllabelwidth}{3em}
\newenvironment{CSLReferences}[2] % #1 hanging-indent, #2 entry-spacing
 {\begin{list}{}{%
  \setlength{\itemindent}{0pt}
  \setlength{\leftmargin}{0pt}
  \setlength{\parsep}{0pt}
  % turn on hanging indent if param 1 is 1
  \ifodd #1
   \setlength{\leftmargin}{\cslhangindent}
   \setlength{\itemindent}{-1\cslhangindent}
  \fi
  % set entry spacing
  \setlength{\itemsep}{#2\baselineskip}}}
 {\end{list}}
\usepackage{calc}
\newcommand{\CSLBlock}[1]{\hfill\break\parbox[t]{\linewidth}{\strut\ignorespaces#1\strut}}
\newcommand{\CSLLeftMargin}[1]{\parbox[t]{\csllabelwidth}{\strut#1\strut}}
\newcommand{\CSLRightInline}[1]{\parbox[t]{\linewidth - \csllabelwidth}{\strut#1\strut}}
\newcommand{\CSLIndent}[1]{\hspace{\cslhangindent}#1}
\setlength{\emergencystretch}{3em} % prevent overfull lines
\providecommand{\tightlist}{%
  \setlength{\itemsep}{0pt}\setlength{\parskip}{0pt}}
\usepackage{booktabs}
\usepackage{longtable}
\usepackage{array}
\usepackage{multirow}
\usepackage{wrapfig}
\usepackage{float}
\usepackage{colortbl}
\usepackage{pdflscape}
\usepackage{tabu}
\usepackage{threeparttable}
\usepackage{threeparttablex}
\usepackage[normalem]{ulem}
\usepackage{makecell}
\usepackage{xcolor}
\usepackage{bookmark}
\IfFileExists{xurl.sty}{\usepackage{xurl}}{} % add URL line breaks if available
\urlstyle{same}
\hypersetup{
  pdftitle={MMRM vs.~Survival Analysis for AD Clinical Trials: A Historical Review and Simulation Study},
  pdfauthor={R.G. Thomas},
  hidelinks,
  pdfcreator={LaTeX via pandoc}}

\title{MMRM vs.~Survival Analysis for AD Clinical Trials: A Historical
Review and Simulation Study}
\author{R.G. Thomas}
\date{2026-03-10}

\begin{document}
\maketitle

{
\setcounter{tocdepth}{2}
\tableofcontents
}
\section{Introduction}\label{introduction}

\subsection{The divergent paths of clinical trial
analysis}\label{the-divergent-paths-of-clinical-trial-analysis}

The statistical analysis of clinical trials has followed markedly
different trajectories across therapeutic areas. In oncology,
time-to-event endpoints analyzed by survival methods remain the
regulatory gold standard, with overall survival serving as the
definitive measure of treatment efficacy (U.S. Food and Drug
Administration 2018a; Fleming and DeMets 1996). In contrast, Alzheimer's
disease (AD) and other neurodegenerative disease trials have converged
on change-from-baseline estimands analyzed by the mixed model for
repeated measures (MMRM) as the primary analytic framework (Donohue and
Aisen 2012; Mallinckrodt, Sanger, et al. 2003). This divergence is not
an accident of history but rather reflects fundamental differences in
disease biology, regulatory philosophy, and the statistical properties
of the available endpoints.

This manuscript traces the historical developments that led to the
current analytic conventions in each field and presents a simulation
study comparing MMRM and survival analysis applied to the same
underlying disease process, using the Clinical Dementia Rating (CDR)
scale as the primary measure.

\subsection{Regulatory origins: the FDA dual-outcome
requirement}\label{regulatory-origins-the-fda-dual-outcome-requirement}

The regulatory framework for AD clinical trials was shaped decisively by
the 1990 FDA guidelines authored by Paul Leber, which required
demonstration of efficacy on two co-primary endpoints: one cognitive
(such as the Alzheimer's Disease Assessment Scale-Cognitive Subscale,
ADAS-Cog) and one global or functional measure (Leber 1990; Rosen et al.
1984). The rationale was that a statistically significant difference on
a psychometric test alone might not reflect clinically meaningful
change; a second, independent measure of global function was needed to
confirm relevance. This dual-outcome requirement established the
precedent that AD trials would be evaluated on the magnitude and
direction of change from baseline, rather than on the occurrence of
discrete clinical events.

The ADAS-Cog, developed by Rosen et al. (1984), and the CDR, codified by
Morris (1993), became the two pillars of AD trial assessment. The CDR in
particular provided both a global staging instrument (the global CDR
score, ranging from 0 to 3) and a sum-of-boxes score (CDR-SB, ranging
from 0 to 18) that offered a quasi-continuous measure of disease
severity suitable for change-score analysis. The 2018 FDA guidance on
early AD further reinforced the change-from-baseline paradigm, while
allowing greater flexibility in endpoint selection for earlier disease
stages (U.S. Food and Drug Administration 2018b).

\subsection{The fall of LOCF and the rise of
MMRM}\label{the-fall-of-locf-and-the-rise-of-mmrm}

Through the 1990s, the standard analytic approach for AD trials was last
observation carried forward (LOCF) combined with analysis of covariance
(ANCOVA) at the final time point. Under LOCF, a participant's last
observed value is imputed for all subsequent missing visits, and the
analysis proceeds as if no data were missing. This approach was
computationally convenient and widely accepted by regulators, but it
rests on the untenable assumption that a participant's trajectory would
have remained flat from the time of dropout---an assumption grossly
violated in progressive diseases where dropout is often driven by
worsening symptoms.

The foundational work of Mallinckrodt et al. (2001a) and Mallinckrodt et
al. (2001b) demonstrated through simulation that LOCF produces
substantially biased treatment effect estimates and inflated Type I
error rates when dropout is differential between treatment arms. In a
progressive disease like AD, where placebo-arm participants tend to drop
out at higher rates due to worsening, LOCF paradoxically imputes
favorable values for the placebo group (freezing them at an earlier,
less impaired state), thereby attenuating the estimated treatment
effect. Conversely, if active-treatment participants drop out due to
adverse events while still relatively well, LOCF preserves their
favorable scores, inflating the apparent benefit.

Mallinckrodt, Sanger, et al. (2003) and Mallinckrodt, Clark, et al.
(2003) extended this work by comparing LOCF to the mixed model for
repeated measures (MMRM), which uses all observed data under a
likelihood-based framework that is valid when data are missing at random
(MAR). Unlike LOCF, MMRM does not impute missing values; instead, it
borrows information from the observed correlation structure across time
points to provide unbiased estimates under MAR. O'Brien et al. (2005)
proposed alternative semi-parametric and non-parametric methods and
found that MMRM performed comparably to their proposed approaches while
making less restrictive assumptions. Molenberghs et al. (2004) provided
the theoretical underpinning for likelihood-based analyses of incomplete
longitudinal data.

By the mid-2000s, the weight of evidence against LOCF was overwhelming,
and the 2010 National Research Council panel report, commissioned by the
FDA, formally recommended against single imputation methods and endorsed
likelihood-based and multiple imputation approaches as preferred
alternatives (National Research Council 2010). The ICH E9(R1) addendum
on estimands further clarified the conceptual framework by requiring
explicit specification of the target estimand---including how
intercurrent events such as treatment discontinuation are
handled---before selecting an analytic method (International Council for
Harmonisation 2019).

\subsection{MMRM as the primary analysis in
AD}\label{mmrm-as-the-primary-analysis-in-ad}

Donohue and Aisen (2012) compared MMRM, slope models, and quadratic
models across five AD Cooperative Study (ADCS) trials and found that
MMRM, which treats time as categorical and imposes no parametric
assumptions on the mean trajectory, provided robust estimates across a
range of disease stages and trial durations. Although continuous time
(slope) models can be more statistically efficient when the linearity
assumption holds, the unstructured mean of MMRM offers protection
against model misspecification in a disease where trajectories are
heterogeneous and potentially nonlinear. Wang et al. (2018) proposed a
disease progression model (DPM) with proportional treatment effects and
demonstrated power gains over MMRM in autosomal dominant AD, though the
DPM requires stronger parametric assumptions that may limit
generalizability.

Today, virtually every pivotal AD trial uses MMRM or a close variant as
the pre-specified primary analysis. The aducanumab EMERGE and ENGAGE
trials (Budd Haeberlein et al. 2022) and the lecanemab Study 201 (Dhadda
et al. 2022) all employed MMRM with change from baseline in CDR-SB as
the primary endpoint, with sensitivity analyses from DPM and other
longitudinal models used to assess robustness.

\subsection{Why survival analysis persists in
oncology}\label{why-survival-analysis-persists-in-oncology}

In oncology, the natural endpoint is death, and overall survival (OS)
has served as the primary measure of treatment efficacy since the
earliest cancer clinical trials. OS is binary, unambiguous, and
universally meaningful. Fleming and DeMets (1996) cautioned against
replacing OS with surrogate endpoints unless rigorous validation
criteria are met, and the FDA has historically required OS data for full
(as opposed to accelerated) approval of cancer drugs (U.S. Food and Drug
Administration 2018a). Even when progression-free survival or response
rate serves as the primary endpoint for accelerated approval,
confirmatory trials typically require OS data.

The persistence of survival analysis in oncology reflects several
features that distinguish it from AD: (1) a clearly defined,
irreversible event (death); (2) a well-understood proportional hazards
framework with established regulatory precedent; (3) a tradition of
interim analyses based on information fractions from event counts; and
(4) the absence of a continuous, validated severity measure that changes
monotonically and is universally accepted across tumor types.

\subsection{Rate of change vs.~time to event in
AD}\label{rate-of-change-vs.-time-to-event-in-ad}

Given that CDR and similar scales can be used to define both change
scores (amenable to MMRM) and threshold crossings (amenable to survival
analysis), the question arises: what is the relative efficiency of these
two approaches when applied to the same underlying data?

Donohue et al. (2011) derived an analytic inflation factor comparing Cox
proportional hazards models of time to a threshold with linear models of
continuous outcomes, and demonstrated through simulation that linear
models consistently yielded greater statistical power for detecting
treatment effects in ADNI data. Li et al. (2019) extended this work to
presymptomatic AD, finding that power was approximately doubled with
repeated continuous measures compared to time-to-progression analysis.
Jönsson et al. (2024) introduced a progression model for repeated
measures (PMRM) that expresses treatment effects in time units (days of
delay) rather than CDR-SB points, offering a clinically interpretable
alternative while maintaining efficiency comparable to MMRM. Ito and
Hutmacher (2014) modeled CDR-SB trajectories to predict time to
clinically meaningful worsening in MCI patients, finding that the median
time to a 1-point change was approximately 2 years---informing both
trial duration and sample size planning.

\subsection{Present study}\label{present-study}

Despite the theoretical and empirical evidence favoring continuous
measures, the comparison between MMRM and survival analysis for AD
endpoints has not been systematically examined in a simulation that uses
the global CDR as both a continuous outcome and the basis for a survival
endpoint defined by conversion from CDR 0.5 (MCI/very mild dementia) to
CDR 1.0 or higher (mild dementia). The present study addresses this gap
through a Monte Carlo simulation that generates realistic CDR
trajectories under a plausible data-generating model and compares the
statistical properties of MMRM and Cox regression applied to the same
simulated trial data.

\section{Methods}\label{methods}

\subsection{Data generating model}\label{data-generating-model}

We simulate a two-arm, placebo-controlled trial with \(n = 200\)
participants per arm, all entering with a baseline global CDR of 0.5.
The primary continuous outcome is the global CDR score observed at
visits scheduled every 6 months over 24 months (5 post-baseline visits).
The survival endpoint is time to conversion from CDR 0.5 to CDR \(\geq\)
1.0.

The latent continuous disease severity \(Y_{ij}^*\) for participant
\(i\) at visit \(j\) is modeled as:

\[
Y_{ij}^* = \beta_0 + \beta_1 t_j + \beta_2 T_i
  + \beta_3 (t_j \times T_i) + b_{0i} + b_{1i} t_j
  + \epsilon_{ij}
\]

where \(t_j\) is visit time in years, \(T_i\) is the treatment indicator
(0 = placebo, 1 = active), \(b_{0i}\) and \(b_{1i}\) are
participant-specific random intercepts and slopes drawn from a bivariate
normal distribution, and \(\epsilon_{ij} \sim N(0, \sigma^2)\) is
residual error.

The observed global CDR is obtained by mapping \(Y^*\) to the CDR
staging categories \{0, 0.5, 1, 2, 3\} using threshold parameters
calibrated to observed CDR distributions in ADNI.

\subsection{Parameter values}\label{parameter-values}

The simulation parameters are selected to reflect plausible progression
rates in a population with MCI due to AD (CDR 0.5 at baseline):

\begin{table}[!h]
\centering
\caption{\label{tab:sim-params}Simulation parameter values.}
\centering
\fontsize{9}{11}\selectfont
\begin{tabular}[t]{lll}
\toprule
Parameter & Value & Description\\
\midrule
$\beta_0$ & 0.5 & Baseline CDR (MCI)\\
$\beta_1$ & 0.30 & Annual rate, placebo (CDR/yr)\\
$\beta_2$ & 0.00 & Baseline treatment effect\\
$\beta_3$ & -0.075 & Treatment x time (25\% slowing)\\
$\sigma_b^2$ (intercept) & 0.04 & Random intercept variance\\
\addlinespace
$\sigma_b^2$ (slope) & 0.02 & Random slope variance\\
$\rho$ & 0.3 & Intercept-slope correlation\\
$\sigma^2$ & 0.01 & Residual variance\\
Dropout rate & 15\%/year & MAR dropout, severity-dependent\\
\bottomrule
\end{tabular}
\end{table}

\subsection{Design specifications}\label{design-specifications}

\begin{itemize}
\tightlist
\item
  \textbf{Sample size:} 200 per arm (400 total)
\item
  \textbf{Visit schedule:} Baseline, 6, 12, 18, 24 months
\item
  \textbf{Survival endpoint:} Time from randomization to first visit
  with global CDR \(\geq\) 1.0
\item
  \textbf{MMRM endpoint:} Change from baseline in latent CDR score at 24
  months
\item
  \textbf{Treatment effect:} 25\% slowing of progression (hazard ratio
  \(\approx\) 0.75 for the survival endpoint)
\end{itemize}

\subsection{Simulation procedure}\label{simulation-procedure}

\begin{Shaded}
\begin{Highlighting}[]
\NormalTok{sim\_trial }\OtherTok{\textless{}{-}} \ControlFlowTok{function}\NormalTok{(}\AttributeTok{n\_per\_arm =} \DecValTok{200}\NormalTok{,}
                      \AttributeTok{n\_visits =} \DecValTok{5}\NormalTok{,}
                      \AttributeTok{beta =} \FunctionTok{c}\NormalTok{(}\FloatTok{0.5}\NormalTok{, }\FloatTok{0.30}\NormalTok{, }\DecValTok{0}\NormalTok{, }\SpecialCharTok{{-}}\FloatTok{0.075}\NormalTok{),}
                      \AttributeTok{sigma\_b =} \FunctionTok{matrix}\NormalTok{(}
                        \FunctionTok{c}\NormalTok{(}\FloatTok{0.04}\NormalTok{, }\FloatTok{0.3} \SpecialCharTok{*} \FunctionTok{sqrt}\NormalTok{(}\FloatTok{0.04} \SpecialCharTok{*} \FloatTok{0.02}\NormalTok{),}
                          \FloatTok{0.3} \SpecialCharTok{*} \FunctionTok{sqrt}\NormalTok{(}\FloatTok{0.04} \SpecialCharTok{*} \FloatTok{0.02}\NormalTok{), }\FloatTok{0.02}\NormalTok{),}
                        \DecValTok{2}\NormalTok{, }\DecValTok{2}
\NormalTok{                      ),}
                      \AttributeTok{sigma\_e =} \FloatTok{0.1}\NormalTok{,}
                      \AttributeTok{dropout\_rate =} \FloatTok{0.15}\NormalTok{) \{}
\NormalTok{  n }\OtherTok{\textless{}{-}} \DecValTok{2} \SpecialCharTok{*}\NormalTok{ n\_per\_arm}
\NormalTok{  trt }\OtherTok{\textless{}{-}} \FunctionTok{rep}\NormalTok{(}\DecValTok{0}\SpecialCharTok{:}\DecValTok{1}\NormalTok{, }\AttributeTok{each =}\NormalTok{ n\_per\_arm)}
\NormalTok{  times }\OtherTok{\textless{}{-}} \FunctionTok{seq}\NormalTok{(}\FloatTok{0.5}\NormalTok{, }\FloatTok{2.5}\NormalTok{, }\AttributeTok{by =} \FloatTok{0.5}\NormalTok{)}

\NormalTok{  re }\OtherTok{\textless{}{-}}\NormalTok{ MASS}\SpecialCharTok{::}\FunctionTok{mvrnorm}\NormalTok{(n, }\AttributeTok{mu =} \FunctionTok{c}\NormalTok{(}\DecValTok{0}\NormalTok{, }\DecValTok{0}\NormalTok{), }\AttributeTok{Sigma =}\NormalTok{ sigma\_b)}

\NormalTok{  records }\OtherTok{\textless{}{-}} \FunctionTok{list}\NormalTok{()}
  \ControlFlowTok{for}\NormalTok{ (i }\ControlFlowTok{in} \FunctionTok{seq\_len}\NormalTok{(n)) \{}
\NormalTok{    y\_star }\OtherTok{\textless{}{-}} \FunctionTok{numeric}\NormalTok{(n\_visits)}
\NormalTok{    dropped }\OtherTok{\textless{}{-}} \ConstantTok{FALSE}
    \ControlFlowTok{for}\NormalTok{ (j }\ControlFlowTok{in} \FunctionTok{seq\_len}\NormalTok{(n\_visits)) \{}
      \ControlFlowTok{if}\NormalTok{ (dropped) }\ControlFlowTok{next}
\NormalTok{      t\_j }\OtherTok{\textless{}{-}}\NormalTok{ times[j]}
\NormalTok{      mu\_ij }\OtherTok{\textless{}{-}}\NormalTok{ beta[}\DecValTok{1}\NormalTok{] }\SpecialCharTok{+}\NormalTok{ beta[}\DecValTok{2}\NormalTok{] }\SpecialCharTok{*}\NormalTok{ t\_j }\SpecialCharTok{+}
\NormalTok{        beta[}\DecValTok{3}\NormalTok{] }\SpecialCharTok{*}\NormalTok{ trt[i] }\SpecialCharTok{+}\NormalTok{ beta[}\DecValTok{4}\NormalTok{] }\SpecialCharTok{*}\NormalTok{ (t\_j }\SpecialCharTok{*}\NormalTok{ trt[i]) }\SpecialCharTok{+}
\NormalTok{        re[i, }\DecValTok{1}\NormalTok{] }\SpecialCharTok{+}\NormalTok{ re[i, }\DecValTok{2}\NormalTok{] }\SpecialCharTok{*}\NormalTok{ t\_j}
\NormalTok{      y\_star[j] }\OtherTok{\textless{}{-}}\NormalTok{ mu\_ij }\SpecialCharTok{+} \FunctionTok{rnorm}\NormalTok{(}\DecValTok{1}\NormalTok{, }\DecValTok{0}\NormalTok{, sigma\_e)}
\NormalTok{      cdr\_obs }\OtherTok{\textless{}{-}} \FunctionTok{cut}\NormalTok{(}
\NormalTok{        y\_star[j],}
        \AttributeTok{breaks =} \FunctionTok{c}\NormalTok{(}\SpecialCharTok{{-}}\ConstantTok{Inf}\NormalTok{, }\FloatTok{0.25}\NormalTok{, }\FloatTok{0.75}\NormalTok{, }\FloatTok{1.5}\NormalTok{, }\FloatTok{2.5}\NormalTok{, }\ConstantTok{Inf}\NormalTok{),}
        \AttributeTok{labels =} \FunctionTok{c}\NormalTok{(}\DecValTok{0}\NormalTok{, }\FloatTok{0.5}\NormalTok{, }\DecValTok{1}\NormalTok{, }\DecValTok{2}\NormalTok{, }\DecValTok{3}\NormalTok{)}
\NormalTok{      )}
\NormalTok{      cdr\_obs }\OtherTok{\textless{}{-}} \FunctionTok{as.numeric}\NormalTok{(}\FunctionTok{as.character}\NormalTok{(cdr\_obs))}
\NormalTok{      p\_drop }\OtherTok{\textless{}{-}} \DecValTok{1} \SpecialCharTok{{-}} \FunctionTok{exp}\NormalTok{(}\SpecialCharTok{{-}}\NormalTok{dropout\_rate }\SpecialCharTok{*} \FloatTok{0.5} \SpecialCharTok{*}
\NormalTok{        (}\DecValTok{1} \SpecialCharTok{+} \FloatTok{0.5} \SpecialCharTok{*}\NormalTok{ (y\_star[j] }\SpecialCharTok{{-}} \FloatTok{0.5}\NormalTok{)))}
      \ControlFlowTok{if}\NormalTok{ (}\FunctionTok{runif}\NormalTok{(}\DecValTok{1}\NormalTok{) }\SpecialCharTok{\textless{}}\NormalTok{ p\_drop) dropped }\OtherTok{\textless{}{-}} \ConstantTok{TRUE}
\NormalTok{      records[[}\FunctionTok{length}\NormalTok{(records) }\SpecialCharTok{+} \DecValTok{1}\NormalTok{]] }\OtherTok{\textless{}{-}} \FunctionTok{data.frame}\NormalTok{(}
        \AttributeTok{id =}\NormalTok{ i, }\AttributeTok{trt =}\NormalTok{ trt[i], }\AttributeTok{visit =}\NormalTok{ j,}
        \AttributeTok{time =}\NormalTok{ t\_j, }\AttributeTok{y\_star =}\NormalTok{ y\_star[j],}
        \AttributeTok{cdr =}\NormalTok{ cdr\_obs}
\NormalTok{      )}
\NormalTok{    \}}
\NormalTok{  \}}
  \FunctionTok{do.call}\NormalTok{(rbind, records)}
\NormalTok{\}}
\end{Highlighting}
\end{Shaded}

For each of \(R = 2000\) replications, we:

\begin{enumerate}
\def\labelenumi{\arabic{enumi}.}
\tightlist
\item
  Generate a complete trial dataset using the model above.
\item
  Apply severity-dependent MAR dropout.
\item
  Fit an MMRM to the change from baseline in the latent CDR score, with
  treatment, visit, treatment-by-visit interaction, and baseline score
  as covariates, using an unstructured covariance matrix.
\item
  Fit a Cox proportional hazards model to time to first CDR \(\geq\)
  1.0, with treatment as the sole covariate.
\item
  Record the treatment effect estimate, standard error, p-value, and
  95\% confidence interval from each model.
\end{enumerate}

\subsection{Analysis plan}\label{analysis-plan}

\begin{Shaded}
\begin{Highlighting}[]
\NormalTok{fit\_mmrm }\OtherTok{\textless{}{-}} \ControlFlowTok{function}\NormalTok{(dat) \{}
\NormalTok{  dat }\OtherTok{\textless{}{-}}\NormalTok{ dat }\SpecialCharTok{|\textgreater{}}
    \FunctionTok{group\_by}\NormalTok{(id) }\SpecialCharTok{|\textgreater{}}
    \FunctionTok{mutate}\NormalTok{(}
      \AttributeTok{baseline =} \FunctionTok{first}\NormalTok{(y\_star),}
      \AttributeTok{change =}\NormalTok{ y\_star }\SpecialCharTok{{-}}\NormalTok{ baseline}
\NormalTok{    ) }\SpecialCharTok{|\textgreater{}}
    \FunctionTok{ungroup}\NormalTok{() }\SpecialCharTok{|\textgreater{}}
    \FunctionTok{filter}\NormalTok{(visit }\SpecialCharTok{\textgreater{}} \DecValTok{0}\NormalTok{) }\SpecialCharTok{|\textgreater{}}
    \FunctionTok{mutate}\NormalTok{(}\AttributeTok{visit\_f =} \FunctionTok{factor}\NormalTok{(visit))}

\NormalTok{  mod }\OtherTok{\textless{}{-}}\NormalTok{ nlme}\SpecialCharTok{::}\FunctionTok{gls}\NormalTok{(}
\NormalTok{    change }\SpecialCharTok{\textasciitilde{}}\NormalTok{ trt }\SpecialCharTok{*}\NormalTok{ visit\_f }\SpecialCharTok{+}\NormalTok{ baseline,}
    \AttributeTok{data =}\NormalTok{ dat,}
    \AttributeTok{correlation =}\NormalTok{ nlme}\SpecialCharTok{::}\FunctionTok{corSymm}\NormalTok{(}\AttributeTok{form =} \SpecialCharTok{\textasciitilde{}}\NormalTok{ visit }\SpecialCharTok{|}\NormalTok{ id),}
    \AttributeTok{weights =}\NormalTok{ nlme}\SpecialCharTok{::}\FunctionTok{varIdent}\NormalTok{(}\AttributeTok{form =} \SpecialCharTok{\textasciitilde{}} \DecValTok{1} \SpecialCharTok{|}\NormalTok{ visit\_f),}
    \AttributeTok{na.action =}\NormalTok{ na.omit}
\NormalTok{  )}
\NormalTok{  trt\_24 }\OtherTok{\textless{}{-}} \FunctionTok{summary}\NormalTok{(mod)}\SpecialCharTok{$}\NormalTok{tTable[}\StringTok{"trt"}\NormalTok{, ]}
  \FunctionTok{c}\NormalTok{(}\AttributeTok{est =}\NormalTok{ trt\_24[}\DecValTok{1}\NormalTok{], }\AttributeTok{se =}\NormalTok{ trt\_24[}\DecValTok{2}\NormalTok{], }\AttributeTok{p =}\NormalTok{ trt\_24[}\DecValTok{4}\NormalTok{])}
\NormalTok{\}}

\NormalTok{fit\_cox }\OtherTok{\textless{}{-}} \ControlFlowTok{function}\NormalTok{(dat) \{}
\NormalTok{  surv\_dat }\OtherTok{\textless{}{-}}\NormalTok{ dat }\SpecialCharTok{|\textgreater{}}
    \FunctionTok{group\_by}\NormalTok{(id) }\SpecialCharTok{|\textgreater{}}
    \FunctionTok{summarise}\NormalTok{(}
      \AttributeTok{trt =} \FunctionTok{first}\NormalTok{(trt),}
      \AttributeTok{event\_time =}\NormalTok{ \{}
\NormalTok{        conv }\OtherTok{\textless{}{-}} \FunctionTok{which}\NormalTok{(cdr }\SpecialCharTok{\textgreater{}=} \DecValTok{1}\NormalTok{)}
        \ControlFlowTok{if}\NormalTok{ (}\FunctionTok{length}\NormalTok{(conv) }\SpecialCharTok{\textgreater{}} \DecValTok{0}\NormalTok{) time[conv[}\DecValTok{1}\NormalTok{]] }\ControlFlowTok{else} \FunctionTok{max}\NormalTok{(time)}
\NormalTok{      \},}
      \AttributeTok{event =} \FunctionTok{as.numeric}\NormalTok{(}\FunctionTok{any}\NormalTok{(cdr }\SpecialCharTok{\textgreater{}=} \DecValTok{1}\NormalTok{)),}
      \AttributeTok{.groups =} \StringTok{"drop"}
\NormalTok{    )}
\NormalTok{  mod }\OtherTok{\textless{}{-}} \FunctionTok{coxph}\NormalTok{(}
    \FunctionTok{Surv}\NormalTok{(event\_time, event) }\SpecialCharTok{\textasciitilde{}}\NormalTok{ trt,}
    \AttributeTok{data =}\NormalTok{ surv\_dat}
\NormalTok{  )}
\NormalTok{  s }\OtherTok{\textless{}{-}} \FunctionTok{summary}\NormalTok{(mod)}
  \FunctionTok{c}\NormalTok{(}
    \AttributeTok{est =} \FunctionTok{coef}\NormalTok{(mod),}
    \AttributeTok{se =}\NormalTok{ s}\SpecialCharTok{$}\NormalTok{coefficients[, }\StringTok{"se(coef)"}\NormalTok{],}
    \AttributeTok{p =}\NormalTok{ s}\SpecialCharTok{$}\NormalTok{coefficients[, }\StringTok{"Pr(\textgreater{}|z|)"}\NormalTok{]}
\NormalTok{  )}
\NormalTok{\}}
\end{Highlighting}
\end{Shaded}

\subsection{Performance metrics}\label{performance-metrics}

For each method, we compute:

\begin{itemize}
\tightlist
\item
  \textbf{Bias:} Mean estimate minus true parameter value
\item
  \textbf{Empirical SE:} Standard deviation of estimates across
  replications
\item
  \textbf{Mean model-based SE:} Mean of estimated standard errors
\item
  \textbf{Power:} Proportion of replications with \(p < 0.05\)
\item
  \textbf{Coverage:} Proportion of 95\% CIs containing the true value
\item
  \textbf{Convergence rate:} Proportion of replications producing valid
  estimates
\end{itemize}

\section{Results}\label{results}

\subsection{Summary statistics}\label{summary-statistics}

\emph{{[}Results table to be populated after simulation execution.{]}}

\subsection{Figures}\label{figures}

\section{Discussion}\label{discussion}

The historical review reveals that the convergence of AD trials on MMRM
and change-score estimands resulted from three reinforcing forces: (1)
the regulatory requirement for continuous cognitive and functional
co-primary endpoints, established by Leber (1990); (2) the decisive
empirical evidence against LOCF as a primary analysis method, led by
Mallinckrodt et al. (2001a) and codified by the National Research
Council (2010) panel; and (3) the statistical efficiency advantages of
modeling repeated continuous data relative to dichotomizing into
time-to-event endpoints, as demonstrated by Donohue et al. (2011) and Li
et al. (2019).

Oncology, by contrast, arrived at a different equilibrium because its
primary endpoint---death---is inherently binary and irreversible,
because the proportional hazards framework has deep regulatory
precedent, and because no single continuous severity measure plays the
role that ADAS-Cog or CDR-SB plays in AD (Fleming and DeMets 1996; U.S.
Food and Drug Administration 2018a).

The simulation results are expected to confirm the theoretical
predictions from the literature: MMRM should yield substantially greater
power than Cox regression when both are applied to CDR data from the
same simulated trial, because the continuous analysis uses information
from the full trajectory of scores at every visit, whereas the survival
analysis reduces this trajectory to a single binary event. The power
advantage should be most pronounced in the MCI-to-mild-dementia
transition modeled here, where conversion rates are modest and many
participants in a 24-month trial will not reach the CDR \(\geq\) 1.0
threshold, resulting in heavy right-censoring that limits the
information available to the Cox model.

Several caveats warrant emphasis. First, the comparison is inherently
asymmetric: the MMRM uses more information from each participant (all
visits) than the Cox model (only the first conversion event). Second,
the survival endpoint is limited by the visit schedule; true conversion
may occur between visits, introducing interval censoring that the
standard Cox model does not accommodate. Third, the relative efficiency
depends on the threshold chosen; a threshold that most participants
cross during follow-up would shift the balance toward the survival
analysis. Fourth, the clinical relevance of the two estimands differs: a
treatment effect expressed as slowing of CDR-SB decline addresses a
different clinical question than one expressed as delay to a diagnostic
milestone. Both may be meaningful to patients and regulators, and the
choice between them should be driven by the estimand framework of
International Council for Harmonisation (2019) rather than by
statistical convenience alone.

\section{Future research}\label{future-research}

\begin{enumerate}
\def\labelenumi{\arabic{enumi}.}
\item
  \textbf{Additional covariance structures.} The present simulation
  assumes MAR dropout. Extensions under MNAR mechanisms (e.g.,
  pattern-mixture models, selection models) would illuminate the
  robustness of each method under more adversarial missingness.
\item
  \textbf{Nonlinear trajectories.} AD progression is not strictly
  linear, particularly in the transition from MCI to dementia.
  Simulation of sigmoidal or piecewise trajectories would test the
  sensitivity of MMRM (which makes no linearity assumption) relative to
  slope-based alternatives.
\item
  \textbf{Alternative survival formulations.} Interval-censored survival
  models and multi-state models that accommodate transitions among CDR
  stages (0.5 \(\to\) 1 \(\to\) 2 \(\to\) 3) could recover some of the
  information lost by standard Cox regression.
\item
  \textbf{Joint modeling.} Joint models of longitudinal CDR trajectories
  and the time-to-conversion event could unify the two analytic
  approaches, potentially offering both the efficiency of longitudinal
  modeling and the clinical interpretability of a time-to-event
  estimand.
\item
  \textbf{Software tools.} An R package implementing the simulation
  framework and providing power calculation utilities for both MMRM and
  survival designs in AD trials would facilitate trial planning.
\item
  \textbf{Applied validation.} Retrospective application of both methods
  to completed trial datasets (e.g., from the ADCS or ADNI) would
  ground-truth the simulation findings in real data with genuine
  missingness patterns and measurement properties.
\item
  \textbf{Estimand alignment.} Further work is needed to clarify how the
  ICH E9(R1) estimand framework should be applied when the same
  underlying data support both change-score and time-to-event estimands,
  particularly regarding the handling of intercurrent events such as
  treatment switching or rescue medication.
\end{enumerate}

\section*{References}\label{references}
\addcontentsline{toc}{section}{References}

\protect\phantomsection\label{refs}
\begin{CSLReferences}{1}{1}
\bibitem[\citeproctext]{ref-buddhaeberlein2022aducanumab}
Budd Haeberlein, S., P. S. Aisen, F. Barkhof, et al. 2022. {``Two
Randomized Phase 3 Studies of Aducanumab in Early {Alzheimer's}
Disease.''} \emph{The Journal of Prevention of Alzheimer's Disease} 9
(2): 197--210. \url{https://doi.org/10.14283/jpad.2022.30}.

\bibitem[\citeproctext]{ref-dhadda2022lecanemab}
Dhadda, Shobha, Michio Kanekiyo, David Li, et al. 2022. {``Consistency
of Efficacy Results Across Various Clinical Measures and Statistical
Methods in the Lecanemab Phase 2 Trial of Early {Alzheimer's}
Disease.''} \emph{Alzheimer's Research \& Therapy} 14 (1): 182.
\url{https://doi.org/10.1186/s13195-022-01129-x}.

\bibitem[\citeproctext]{ref-donohue2012mmrm}
Donohue, M. C., and P. S. Aisen. 2012. {``Mixed Model of Repeated
Measures Versus Slope Models in {Alzheimer's} Disease Clinical
Trials.''} \emph{The Journal of Nutrition, Health \& Aging} 16 (4):
360--64. \url{https://doi.org/10.1007/s12603-012-0047-7}.

\bibitem[\citeproctext]{ref-donohue2011relative}
Donohue, M. C., A. C. Gamst, R. G. Thomas, et al. 2011. {``The Relative
Efficiency of Time-to-Threshold and Rate of Change in Longitudinal
Data.''} \emph{Contemporary Clinical Trials} 32 (5): 685--93.
\url{https://doi.org/10.1016/j.cct.2011.04.007}.

\bibitem[\citeproctext]{ref-fleming1996surrogate}
Fleming, Thomas R., and David L. DeMets. 1996. {``Surrogate End Points
in Clinical Trials: Are We Being Misled?''} \emph{Annals of Internal
Medicine} 125 (7): 605--13.
\url{https://doi.org/10.7326/0003-4819-125-7-199610010-00011}.

\bibitem[\citeproctext]{ref-ich2019e9r1}
International Council for Harmonisation. 2019. \emph{{ICH E9(R1)}:
Addendum on Estimands and Sensitivity Analysis in Clinical Trials to the
Guideline on Statistical Principles for Clinical Trials}. ICH.

\bibitem[\citeproctext]{ref-ito2014predicting}
Ito, Kaori, and Matthew M. Hutmacher. 2014. {``Predicting the Time to
Clinically Worsening in Mild Cognitive Impairment Patients and Its
Utility in Clinical Trial Design by Modeling a Longitudinal {Clinical
Dementia Rating Sum of Boxes} from the {ADNI} Database.''} \emph{Journal
of Alzheimer's Disease} 40 (4): 967--79.
\url{https://doi.org/10.3233/JAD-132090}.

\bibitem[\citeproctext]{ref-jonsson2024progression}
Jönsson, Linus, Milana Ivkovic, Alireza Atri, et al. 2024.
{``Progression Analysis Versus Traditional Methods to Quantify Slowing
of Disease Progression in {Alzheimer's} Disease.''} \emph{Alzheimer's
Research \& Therapy} 16 (1): 48.
\url{https://doi.org/10.1186/s13195-024-01413-y}.

\bibitem[\citeproctext]{ref-leber1990guidelines}
Leber, P. 1990. \emph{Guidelines for the Clinical Evaluation of
Antidementia Drugs, First Draft}. U.S. Food; Drug Administration.

\bibitem[\citeproctext]{ref-li2019relative}
Li, Dan, Samuel Iddi, Paul S. Aisen, Wesley K. Thompson, and Michael C.
Donohue. 2019. {``The Relative Efficiency of Time-to-Progression and
Continuous Measures of Cognition in Presymptomatic {Alzheimer's}
Disease.''} \emph{Alzheimer's \& Dementia: Translational Research \&
Clinical Interventions} 5: 308--18.
\url{https://doi.org/10.1016/j.trci.2019.04.004}.

\bibitem[\citeproctext]{ref-mallinckrodt2003response}
Mallinckrodt, Craig H., W. Scott Clark, Raymond J. Carroll, and Geert
Molenberghs. 2003. {``Assessing Response Profiles from Incomplete
Longitudinal Clinical Trial Data Under Regulatory Considerations.''}
\emph{Journal of Biopharmaceutical Statistics} 13 (2): 179--90.
\url{https://doi.org/10.1081/BIP-120019265}.

\bibitem[\citeproctext]{ref-mallinckrodt2001dropout}
Mallinckrodt, Craig H., W. Scott Clark, and Stacy R. David. 2001a.
{``Accounting for Dropout Bias Using Mixed-Effects Models.''}
\emph{Journal of Biopharmaceutical Statistics} 11 (1--2): 9--21.
\url{https://doi.org/10.1081/BIP-100104194}.

\bibitem[\citeproctext]{ref-mallinckrodt2001typei}
Mallinckrodt, Craig H., W. Scott Clark, and Stacy R. David. 2001b.
{``Type {I} Error Rates from Mixed Effects Model Repeated Measures
Versus Fixed Effects {ANOVA} with Missing Values Imputed via Last
Observation Carried Forward.''} \emph{Drug Information Journal} 35 (4):
1215--25. \url{https://doi.org/10.1177/009286150103500418}.

\bibitem[\citeproctext]{ref-mallinckrodt2003assessing}
Mallinckrodt, Craig H., Todd M. Sanger, Sanjay Dubé, et al. 2003.
{``Assessing and Interpreting Treatment Effects in Longitudinal Clinical
Trials with Missing Data.''} \emph{Biological Psychiatry} 53 (8):
754--60. \url{https://doi.org/10.1016/S0006-3223(02)01867-X}.

\bibitem[\citeproctext]{ref-molenberghs2004analyzing}
Molenberghs, Geert, Herbert Thijs, Ivy Jansen, et al. 2004. {``Analyzing
Incomplete Longitudinal Clinical Trial Data.''} \emph{Biostatistics} 5
(3): 445--64. \url{https://doi.org/10.1093/biostatistics/kxh001}.

\bibitem[\citeproctext]{ref-morris1993cdr}
Morris, John C. 1993. {``The {Clinical Dementia Rating} ({CDR}): Current
Version and Scoring Rules.''} \emph{Neurology} 43 (11): 2412--14.
\url{https://doi.org/10.1212/WNL.43.11.2412-a}.

\bibitem[\citeproctext]{ref-nrc2010missing}
National Research Council. 2010. \emph{The Prevention and Treatment of
Missing Data in Clinical Trials}. The National Academies Press.
\url{https://doi.org/10.17226/12955}.

\bibitem[\citeproctext]{ref-obrien2005semiparametric}
O'Brien, Peter C., David Zhang, and Kent R. Bailey. 2005.
{``Semi-Parametric and Non-Parametric Methods for Clinical Trials with
Incomplete Data.''} \emph{Statistics in Medicine} 24 (3): 341--58.
\url{https://doi.org/10.1002/sim.1963}.

\bibitem[\citeproctext]{ref-rosen1984adas}
Rosen, W. G., R. C. Mohs, and K. L. Davis. 1984. {``A New Rating Scale
for {Alzheimer's} Disease.''} \emph{The American Journal of Psychiatry}
141 (11): 1356--64. \url{https://doi.org/10.1176/ajp.141.11.1356}.

\bibitem[\citeproctext]{ref-fda2018oncology}
U.S. Food and Drug Administration. 2018a. \emph{Clinical Trial Endpoints
for the Approval of Cancer Drugs and Biologics. Guidance for Industry}.
Center for Drug Evaluation; Research.

\bibitem[\citeproctext]{ref-fda2018earlyad}
U.S. Food and Drug Administration. 2018b. \emph{Early {Alzheimer's}
Disease: Developing Drugs for Treatment. Guidance for Industry}. Center
for Drug Evaluation; Research.

\bibitem[\citeproctext]{ref-wang2018dpm}
Wang, Guoqiao, Scott Berry, Chengjie Xiong, et al. 2018. {``A Novel
Cognitive Disease Progression Model for Clinical Trials in
Autosomal-Dominant {Alzheimer's} Disease.''} \emph{Statistics in
Medicine} 37 (21): 3047--55. \url{https://doi.org/10.1002/sim.7811}.

\end{CSLReferences}

\end{document}
